\chapter{Simulación Circuital en LTspice}
\label{chap:simulacion}

Con el propósito de validar computacionalmente el diseño analítico y contrastar el comportamiento del circuito,
se implementó y simuló la etapa amplificadora en el software LTspice. Se estudió el punto estático de
polarización en continua (.op) y la respuesta dinámica en el dominio del tiempo (.tran) ante excitaciones senoidales.

\section{Directivas de Simulación y Modelos Empleados}

Para reproducir de forma realista el transistor utilizado en el laboratorio, se adoptó el modelo SPICE del BC546C/BC547
con ganancia directa ajustada a $B_f = 284$, idéntica al valor medido en el multímetro de banco.

En el listado~\ref{list:directivas_spice_tp4} se transcriben las directivas de simulación y el modelo configurado en
el archivo esquemático \texttt{ValoresComerciales.asc}.

\begin{lstlisting}[style=ltspice, caption={Directivas de simulación y modelo SPICE del transistor en LTspice.}, label={list:directivas_spice_tp4}]
* Modelo SPICE del transistor NPN (Philips) con Bf = 284
.model BC546C NPN(Is=21.43f Xti=3 Eg=1.11 Vaf=74.03 Bf=284
+ Ise=43.11f Ne=1.442 Ikf=97.09m Nk=.5 Xtb=1.5 Br=9.837
+ Isc=144.3f Nc=1.522 Ikr=2.476 Rc=1 Cjc=5.25p Mjc=.3147
+ Vjc=.5697 Fc=.5 Cje=11.5p Mje=.6115 Vje=.8194 Tr=10n
+ Tf=443.2p Itf=.7 Vtf=4 Xtf=4 Vceo=65 Icrating=100m mfg=Philips)

* Fuente de senal senoidal de entrada: Amplitud 10 mVp (20 mVpp), f = 1 kHz
VBB n002 0 SINE(0 10m 1k)

* Directiva de analisis en el dominio del tiempo
.tran 50m
\end{lstlisting}

La directiva \lstinline[style=ltspice]{.tran 50m} ejecuta un análisis temporal durante un intervalo total de $50\ms$
(correspondiente a 50 ciclos completos de la señal de audio a $1\kHz$). Este tiempo de integración permite superar
holgadamente los transitorios iniciales de carga de los capacitores de acoplo ($C_1 = C_2 = 22\uF$), alcanzando
el régimen permanente periódico estable.

\section{Simulación del Punto de Operación en Continua (.op)}

La simulación del punto de trabajo estático arrojó los siguientes valores numéricos volcados en \texttt{op.txt}:
\begin{itemize}
  \item Tensión continua de alimentación: $V(v_c) = V_{CC} = 5.000\V$.
  \item Tensión de base respecto a masa: $V(b) = V_B = 1.683\V$.
  \item Tensión de emisor respecto a masa: $V(v_{re}) = V_E = 1.048\V$.
  \item Tensión continua colector-emisor: $V_{CEQ} = V(v_c) - V(v_{re}) = 5.000\V - 1.048\V = 3.952\V$.
  \item Tensión base-emisor en conducción: $V_{BE} = V(b) - V(v_{re}) = 1.683\V - 1.048\V = 0.635\V$.
  \item Corriente de colector de reposo: $I_{CQ} = 1.0429\mA \approx 1.043\mA$.
  \item Corriente de emisor de reposo: $I_{EQ} = 1.0476\mA \approx 1.048\mA$.
  \item Corriente de base de reposo: $I_{BQ} = 4.637\uA$.
  \item Corrientes en el divisor: $I(R_1) = 35.81\uA$ e $I(R_2) = 40.45\uA$.
\end{itemize}

Comparando con el cálculo analítico normalizado ($I_{CQ} = 1.012\mA$ y $V_{CEQ} = 3.985\V$), la coincidencia es
notable, presentando una discrepancia de tan solo $+3.06\%$ en corriente y $-0.83\%$ en tensión. Esta mínima diferencia
proviene de la caída base-emisor no lineal del modelo SPICE ($V_{BE} = 0.635\V$ en lugar de los $0.70\V$ supuestos
idealmente).

\section{Análisis de Formas de Onda y Ganancia Dinámica}

Se extrajeron los datos
del régimen permanente periódico entre $48\ms$ y $50\ms$. En la figura~\ref{fig:sim_waveforms} se presentan las señales
superpuestas de entrada ($v_i$) y salida sobre la carga ($v_L$).

\begin{figure}[!ht]
  \centering
  \begin{tikzpicture}
  \begin{axis}[
    width=0.92\textwidth,
    height=7.2cm,
    xlabel={Tiempo $[\mathrm{ms}]$},
    ylabel={Tensión $[\mathrm{mV}]$},
    xmin=48.0, xmax=50.0,
    ymin=-12, ymax=12,
    grid=both,
    grid style={line width=.1pt, draw=gray!30},
    major grid style={line width=.2pt,draw=gray!50},
    legend pos=north east,
    legend cell align={left},
    every axis plot/.append style={very thick}
  ]
    \addplot[color=blue, smooth] table [x=time, y=vi] {figures/waveforms_data.dat};
    \addlegendentry{Entrada $v_i(t)$ ($20.00\mVpp$)}

    \addplot[color=green!60!black, dashed, smooth] table [x=time, y=vl] {figures/waveforms_data.dat};
    \addlegendentry{Salida sobre carga $v_L(t)$ ($19.47\mVpp$, $A_v \approx 0.974$)}
  \end{axis}
\end{tikzpicture}

  \caption{Formas de onda en régimen permanente a $1\kHz$: respuesta temporal de entrada $v_i(t)$ y salida $v_L(t)$ sobre $R_L$.}
  \label{fig:sim_waveforms}
\end{figure}

A partir de los datos crudos simulados se constata:
\begin{itemize}
  \item Tensión de entrada pico a pico: $v_i = 20.00\mVpp$ (amplitud pico de $10.00\mV$).
  \item Tensión de salida sobre la carga: $v_L = 19.474\mVpp$ (amplitud pico de $9.737\mV$).
\end{itemize}

La ganancia de tensión simulada resulta:
\begin{equation}
  A_{v,\mathrm{sim}} = \frac{v_L}{v_i} = \frac{19.474\mVpp}{20.000\mVpp} \approx 0.9737
  \label{eq:Av_simulada}
\end{equation}

Asimismo, en la figura~\ref{fig:sim_waveforms} se verifica visualmente que ambas ondas se encuentran rigurosamente en
fase ($0^\circ$), confirmando el comportamiento intrínseco de seguidor emisor sin inversión de polaridad.
